\begin{frame}[t,fragile]{べき乗法の収束}
  \begin{itemize}
    %\setlength{\itemsep}{1em}
  \item 実装上の注意
  \begin{itemize}
    \item 反復の途中でベクトルの大きさが大きくなりすぎるとオーバーフローするので、各ステップで正規化する
      \[
      v_{k+1} = \frac{A v_k}{|A v_k|}
      \]
  \end{itemize}
  \item べき乗法による絶対値最大の固有値の近似値(Rayleigh商)
    \[
    \mu_k = \frac{v_{k}^T A v_{k}}{v_{k}^T v_k} = \lambda_1 \Big[ 1 + O\Big( \big|\frac{\lambda_2}{\lambda_1} \big|^{2k} \Big) \Big]
    \]
  \item 誤差の収束
    \[
    \mu_k  \approx \lambda_1 + e^{-2k \ln (|\lambda_1/\lambda_2|)}
    \]
  \item $1 / \ln (|\lambda_1/\lambda_2|)$ 程度の反復が必要
  \item $|\lambda_1|$と$|\lambda_2|$が近い場合には、反復回数が非常に多くなる
  \end{itemize}
\end{frame}
